%&latex
\documentclass[11pt]{amsart}
%%\usepackage[T1]{fontenc}
\usepackage[applemac]{inputenc}
\usepackage[american, frenchb]{babel}

\def\ladate{14 Aožt 2002}
\def\from{From}

\newif\ifarxiv
\arxivtrue
%%\swapnumbers

\newtheorem{theoreme}{ThŽorme}[]
\newtheorem{corollaire}[theoreme]{Corollaire}
\newtheorem{lemme}[theoreme]{Lemme}

%\newtheorem{lemma}[theoreme]{Lemma}
%\newtheorem{theorem}[theoreme]{Theorem}
%\newtheorem{corollary}[theoreme]{Corollary}

\theoremstyle{definition}
\newtheorem{definition}{Definition}
\newtheorem{note}{Note}
\newtheorem{remarque}{Remarque}

%\newtheorem{remark}[theoreme]{Remark}


%%\numberwithin{equation}{section}

\newcommand\Un{{\mathbf 1}}
\newcommand\CC{{\mathbb C}}
\newcommand\NN{{\mathbb N}}
\newcommand\QQ{{\mathbb Q}}
\newcommand\RR{{\mathbb R}}
\newcommand\ZZ{{\mathbb Z}}
\newcommand\HH{{\mathbb H}}


\newcommand\cA{{\mathcal A}}
\newcommand\cB{{\mathcal B}}
\newcommand\cE{{\mathcal E}}
\newcommand\cF{{\mathcal F}}
\newcommand\cG{{\mathcal G}}
\newcommand\cS{{\mathcal S}}
\newcommand\cC{{\mathcal C}}
\newcommand\cU{{\mathcal U}}
\newcommand\cV{{\mathcal V}}
\newcommand\cK{{\mathcal K}}
\newcommand\cD{{\mathcal D}}
\newcommand\cL{{\mathcal L}}

\let\wh=\widehat
\let\wt=\widetilde

\renewcommand{\Re}{{\rm Re}}
\renewcommand{\Im}{{\rm Im}}

\let\la=\lambda
%%\selectlanguage{french}

\begin{document}

\ifarxiv\rightline{\ladate}\else
\textbf{Analyse Harmonique/\textit{Harmonic Analysis}}\hfill\ladate\break\fi
\ifarxiv\rightline{math/0208121}\fi

\bigskip

\title[Sur les Espaces de Sonine]{Sur les ``Espaces de Sonine''
associ\'es par de~Branges \`a la transformation de Fourier}

\author{Jean-Fran\c cois Burnol}


\ifarxiv\else\thanks{\textbf{Note prŽsentŽe par J.-P. Kahane}}\fi


\begin{abstract}
Nous avons obtenu des formules explicites reprŽsentant les
fonctions $E(z)$ apparaissant dans la thŽorie des  ``Espaces de Sonine''
associŽs par De~Branges ˆ la transformation de Fourier. 
\end{abstract}

\maketitle

\selectlanguage{american}
\title{\textbf{On the ``Sonine Spaces'' associated by de~Branges
to the Fourier Transform}}

\begin{abstract}
We have obtained explicit formulae representing the functions $E(z)$
arising in the theory of the ``Sonine Spaces'' associated by De~Branges
to the Fourier Transform.
\end{abstract}

\maketitle

\selectlanguage{french}
\bigskip

\begin{small}
Jean-Fran\c cois Burnol, UniversitŽ Lille 1, UFR de MathŽmatiques, 
CitŽ Scientifique M2, F-59655 Villeneuve d'Ascq Cedex, France\\
\texttt{burnol@agat.univ-lille1.fr}
\end{small}

\parindent=0pt

\baselineskip = 16pt
\parskip=\the\baselineskip

%\section{Abridged English Version}
%
\centerline{\hbox to4cm{\leaders \hrule height1pt \hfill}}
%
\selectlanguage{french}

\section{Espaces et Fonctions de de~Branges}

Nous commencerons par rappeler quelques ŽlŽments de base de la thŽorie
de de~Branges des ``Espaces de Hilbert de fonctions entires''
\cite{bra}. Tout ce dont nous aurons besoin est aussi exposŽ et dŽmontrŽ dans les
pages 220 ˆ 228 du livre \cite{dymkean2} de Dym et McKean. Soit $L$
une droite dans le plan complexe, bordant un demi-plan ouvert $L_+$
qui est toujours chez de~Branges le demi-plan supŽrieur et qui sera
toujours chez nous le demi-plan $\Re(s) > 1/2$. Soit $w^\#$ le
symŽtrique orthogonal du nombre complexe $w$ par rapport ˆ $L$, et
notons $F^\#(w) = \overline{F(w^\#)}$. Pour tout espace non nul de
Hilbert de fonctions entires tel que \textit{(a)} l'Žvaluation en
chaque $w$ est continue, \textit{(b)} $F\mapsto F^\#$ est une
isomŽtrie (anti-unitaire), et \textit{(c)} si $F(\gamma) = 0$ alors
$G(w) = (w-\gamma^\#)/(w-\gamma) F(w)$ appartient ˆ l'espace et est de
mme norme que $F$, il existe par un thŽorme de de~Branges une (non
unique) fonction entire $E(w)$  telle que $|E(w)| > |E(w^\#)|$ pour
tout $w\in L_+$ et telle que l'espace de fonctions entires considŽrŽ
co•ncide isomŽtriquement avec l'espace $\cB(E)$ dŽfini ainsi: $F$ est
dans $\cB(E)$ si $F(w)/E(w)$ est dans l'espace de Hardy $\HH^2(L_+)$
et si $F^\#(w)/E(w)$ l'est aussi. La norme hilbertienne est obtenu par
la premire inclusion dans $\HH^2(L_+)$. Il faut prŽciser quel
multiple de la mesure de Lebesgue est pris sur la droite $L$: chez
de~Branges $L = \RR$ avec la mesure de Lebesgue, chez nous $L =
1/2 + i \RR$ avec $1/2\pi$ fois la mesure de Lebesgue. Un autre
thŽorme de de~Branges nous informe que $\cB(E)$ dŽfini ainsi est
toujours non nul et est un espace de Hilbert, pour toute fonction
entire $E(w)$ satisfaisant ˆ la condition ci-dessus. 
\let\ol = \overline

De plus on sait calculer exactement en fonction de $E$ les produits
scalaires entre Žvaluateurs. DorŽnavant $L_+ = \{\Re(s)>1/2\}$. Alors
$w^\# = 1 - \ol{w}$. On supposera que la ``condition de rŽalitŽ'' est
satisfaite, c'est-ˆ-dire que pour tout $F(w)$ dans l'espace $\overline{F(\overline{w})}$
est dans l'espace avec la mme norme.  On peut alors choisir
$E(w)$ de sorte que $\ol{E(w)} = E(\ol{w})$. On Žcrit $E = A - iB$
avec $A$ ``paire'' ($A(1-w) = A(w)$) et $B$ ``impaire''. Les fonctions
$A$ et $B$ ont tous leurs zŽros sur la droite critique. Les autres
fonctions $E$ dŽfinissant le mme espace (avec la mme norme) et
rŽelles sur l'axe rŽel sont  connues exactement: il s'agit  des
fonctions $kA - iB/k$ pour $k\in \RR, k\neq 0$. 
%%
%Lorsque la fonction
%$A$ ne s'annule pas en $1/2$ on peut spŽcifier canoniquement un unique
%$E$ (fonction rŽelle sur l'axe rŽel) en lui imposant la condition
%$E(1/2) = 1$. Mais d'autres normalisations peuvent tre plus
%appropriŽes.
%%
La fonction (la condition de rŽalitŽ est ici supposŽe vŽrifiŽe)
%%
\begin{equation}
K(z_1, z_2) = \frac{E(z_1)E(z_2) - E(1-z_1)E(1-z_2)}{z_1 + z_2 -
1}
\end{equation}
%%
\begin{equation}\label{noyau}
K(z_1, z_2) 
= 2\frac{(-iB(z_1))A(z_2)+A(z_1)(-iB(z_2))}{z_1 + z_2 - 1}
\end{equation}
%%
est comme fonction de $z_2$ Žgale ˆ l'Žvaluateur en $\ol{z_1}$ et
comme fonction de $z_1$ Žgale ˆ l'Žvaluateur en $\ol{z_2}$.

\section{Espaces de Sonine}

\let\la=\lambda

Soit $\la>0$. Soit $\cF$ la transformation de Fourier
($\cF(\phi)(x) = \int_\RR \exp(2\pi ixy)\phi(y)dy$). Il est connu
classiquement (voir \cite[sect. 2.9]{dymkean1}) que $L^2(-\la,\la) +
\cF L^2(-\la,\la)$ est un sous espace \emph{fermŽ} de l'espace de
Hilbert $L^2(\RR)$ (cela se dŽduit aisŽment en particulier du
caractre compact de l'opŽrateur (ˆ noyau) $F_\la = P_\la \cF P_\la$, o
$P_\la$ est la projection orthogonale sur $L^2(-\la,\la)$). Bien sžr
il s'agit d'un sous-espace propre et son complŽment perpendiculaire
est donc non-nul: il s'agit de l'espace de Hilbert des fonctions de
carrŽs intŽgrables, nulles et de Fourier nulles dans $(-\la,\la)$.
Suivant de~Branges \cite{bra}, nous dirons qu'une fonction est une
fonction de Sonine si elle satisfait ˆ cette condition. 

L'argument le plus simple pour montrer concrtement l'existence de
fonctions de Sonine nous a ŽtŽ communiquŽ par le Professeur Kahane
\footnote{Lettre adressŽe ˆ l'auteur, Mars 2002}: en appliquant un
polyn™me appropriŽ en $x$ et $d/dx$ ˆ la distribution de Poisson
$\sum_{n\in\ZZ} \delta(x-n)$, on obtient une distribution tempŽrŽe
nulle et de Fourier nulle dans $(-A,A)$ avec $A$ arbitraire. Aprs
rŽgularisation par multiplication et convolution additive par des
fonctions tests appropriŽes on obtient une fonction (non-nulle\dots)
dans la classe de Schwartz (donc de carrŽ intŽgrable) ayant la
propriŽtŽ de Sonine, avec $\la>0$ arbitrairement grand.

Dans \cite{jfcopoisson} nous montrons que la
\emph{convolution multiplicative} permet d'associer ˆ toute distribution ayant
la propriŽtŽ d'annulation de Sonine une formule de \emph{co-Poisson} portant sur
des fonctions de Sonine, et
qui gŽnŽralise les formules de co-Poisson associŽes dans \cite{jfhab} ˆ
la fonction dzta de Riemann et aux sŽries $L$ de
Hecke-Tate. Il appara”t d'ailleurs fort utile pour l'Žtude des
fonctions de carrŽs intŽgrables de de~Branges-Sonine de disposer de
rŽsultats gŽnŽraux sur les distributions ayant la mme propriŽtŽ
d'annulation (propriŽtŽ qu'il convient aussi de gŽnŽraliser, par exemple en
remplaant la condition ``nulle dans $(-\la,\la)$'' par la condition ``polynomiale dans
$(-\la,\la)$ de degrŽ $\leq M$''). Dans la suite nous nous restreignons aux fonctions et
distributions \emph{paires}\footnote{Quelques ajustements simples partout dans ce qui suit
comme remplacer les cosinus par des sinus donnent la thŽorie pour le cas impair},
pour lesquelles $\cF$ devient la
transformŽe en cosinus $\cF_+(\phi)(x) = 2\int_0^\infty \cos(2\pi
xy)\phi(y)dy$. Nous dŽsignons par $K_\la$ l'espace des fonctions de
Sonine paires et de carrŽs intŽgrables. Nous utiliserons le rŽsultat
suivant:

%\begin[{\cite{jfcopoisson}}]{theoreme}
\begin{theoreme}[{\cite{jfcopoisson}}]
Soit $D$ une distribution tempŽrŽe paire qui s'annule sur
$(-\la,\la)$. On (rappelle que l'on) peut donner un sens ˆ $\wh{D}(s) = \int_0^\infty
D(t)t^{-s}dt$ comme fonction analytique de $s$ pour $\Re(s)\gg 1$. Si
$\cF_+(D)$ est ˆ nouveau nulle dans $(-\la,\la)$ alors $\wh{D}(s)$ est
une fonction entire de $s$, qui a des zŽros triviaux en $-2n$,
$n\in\NN$. De plus on a l'Žquation fonctionnelle 
$$\pi^{-s/2}\Gamma(s/2)\wh{D}(s) =
\pi^{-(1-s)/2}\Gamma((1-s)/2)\wh{\cF_+(D)}(1-s)$$
\end{theoreme}

ConsidŽrons alors l'espace de Hilbert non-nul $\cS_\la$ des fonctions
entires $F(w) = \pi^{-w/2}\Gamma(w/2)\wh{f}(w)$, o $f$ parcourt
$K_\la$. Il est aisŽ de vŽrifier que $\cS_\la$ vŽrifie les axiomes
\textit{(a)}, \textit{(b)}, et \textit{(c)}, et est donc un espace de
de~Branges, par rapport ˆ la droite critique, et qui vŽrifie la
condition de rŽalitŽ (De~Branges \cite{bra64, bra} a associŽ plus
gŽnŽralement de tels espaces aux transformations de Hankel de
paramtre $\nu$, $\nu>-1$). Le problme (qui est Žquivalent ˆ celui du
calcul des produits scalaires entre Žvaluateurs dans $K_\la$) Žtait donc posŽ de dŽterminer
les fonctions entires $\cA_\la(w)$, $\cB_\la(w)$ et $\cE_\la(w) =
\cA_\la(w) - i \cB_\la(w)$ qui permettent d'Žcrire l'identitŽ
(y-compris pour la norme)  $\cS_\la = \cB(\cE_\la)$. Nous donnons dans
la prŽsente Note une solution explicite ˆ ce problme.

\section{Un problme de projection orthogonale}

Soit
$P_\la$ la projection orthogonale sur $L^2(-\la,\la)$ et $\wt{P_\la}=
\cF P_\la \cF^* = \cF^* P_\la \cF$ la projection sur
$\cF(L^2(-\la,\la))$. Soit $L^2(\RR)^{\rm pair}$ l'espace des fonctions paires 
avec comme norme $\|f\|^2=\int_0^\infty |f(t)|^2dt$. 
Sur $L^2(\RR)^{\rm pair}$ on a  $\wt{P_\la} =
\cF_+ P_\la \cF_+$. Soit 
$K_\la$ son sous-espace de Sonine, et $G_\la =
P_\la(L^2(\RR)^{\rm pair}) + \wt{P_\la}(L^2(\RR)^{\rm pair})$. 

\begin{lemme}
L'espace de Sonine $K_\la$ est le noyau de l'opŽrateur $P_\la + \wt{P_\la}$ (sur 
$L^2(\RR)^{\rm pair}$). La 
restriction de cet opŽrateur ˆ $G_\la$ est 
auto-adjoint et de Fredholm (i.e. de la forme 1+compact) et est inversible. La
projection orthogonale de $f\in L^2(\RR)^{\rm pair}$ sur $G_\la$ est
$(P_\la + \wt{P_\la})^{-1} (P_\la(f) + \wt{P_\la}(f))$.
\end{lemme}

Si $P_\la(f) + \wt{P_\la}(f) = 0$ alors $g=P_\la(f)=-\wt{P_\la}(f) = 0$
 et $f$ est dans $K_\la$.
Donc si $k\in G_\la$ a la mme image sous $P_\la + \wt{P_\la}$ que $f$
alors $f-k$ est dans $K_\la$ et $k$ est la projection orthogonale de $f$ sur $G_\la$.
ConsidŽrons $\phi: L^2(-\la,\la)^{\rm pair}\oplus
L^2(-\la,\la)^{\rm pair} \to G_\la$ qui envoie $(u,v)$ sur $u +
\cF_+(v)$. L'opŽrateur bornŽ $\phi$ est injectif et surjectif donc un
isomorphisme topologique qui conjugue $P_\la + \wt{P_\la}$ ˆ
l'opŽrateur (lui-aussi auto-adjoint) $(u,v) \mapsto (u + F_\la v,
F_\la u + v)$, avec $F_\la = P_\la \cF_+ P_\la$ compact auto-adjoint
de norme strictement infŽrieure ˆ $1$. L'inverse de cet opŽrateur est
$(u,v) \mapsto ((1 - D_\la)^{-1} (u - F_\la v), (1 - D_\la)^{-1} (-
F_\la u + v))$. Dans cette formule  $D_\la = F_\la^2$  peut-tre
remplacŽ par le noyau $P_\la \cF^* P_\la \cF P_\la$ puisqu'il n'agit que sur
des fonctions paires. On note que $F_\la$ a les mmes vecteurs propres
que $D_\la$: ce sont les ``fonctions prolates sphŽroidales'' (paires)
$e_{2n}$, $n\geq0$ (\cite{slepian}). Ainsi:

\begin{lemme} 
Les vecteurs propres de l'opŽrateur (dans $G_\la$) 
de Fredholm $P_\la
+ \wt{P_\la}$ sont les vecteurs $e_{2n} \pm \cF_+(e_{2n})$, $n\in\NN$. 
Ils forment une base orthogonale de $G_\la$ puisque l'opŽrateur est
auto-adjoint.
\end{lemme}

Il est donc possible d'Žcrire la projection orthogonale sur l'espace
de Sonine en utilisant la base orthogonale de $G_\la$ ci-dessus, mais
nous ne ferons pas usage de telles formules.


\begin{remarque} 
L'auteur doit ˆ M.~Balazard et ƒ.~Saias (communication privŽe,
dŽcembre 2001) l'observation que les fonctions sphŽroidales per\-mettent
d'Žcrire une base orthogonale de $G_\la$ (leur base orthogonale
diffre lŽgrement de celle apparaissant ici).
\end{remarque}

Si l'on combine les formules prŽcŽdentes on obtient:

\begin{theoreme}
Soit $f(x)$ une fonction paire de carrŽ intŽgrable. Sa projection
orthogonale sur l'espace de Sonine $K_\la$ est donnŽe par la formule:
%%
$$\pi_\la(f) = f - (1 - D_\la)^{-1}\Big(P_\la(f) - F_\la\cF_+(f)\Big)
- \cF_+(1 - D_\la)^{-1}\Big(P_\la\cF_+(f) -  F_\la(f)\Big)$$
%%
On a $F_\la = P_\la \cF_+ P_\la$ et $D_\la = F_\la^2$ 
agissant sur $L^2(-\la,\la)^{\rm pair}$. On peut remplacer $D_\la$ par le noyau 
$\sin(2\pi\la(x-y))/\pi(x-y)$ restreint ˆ $L^2(-\la,\la)^{\rm pair}$.
\end{theoreme}

\begin{corollaire}
Si $f|_{(-\la,\la)} = 0$ alors 
%%
$$\pi_\la(f) = f\Big|_{|t|>\la} - \Big(\cF_+(1 -
D_\la)^{-1}P_\la\cF_+(f)\Big)\Big|_{|t|>\la}$$
%%
Si $f = \cF_+(f)$ alors
%%
$$\pi_\la(f) = f - (1 + \cF_+)(1+F_\la)^{-1}P_\la(f)$$
%%
Si $f = -\cF_+(f)$ alors
$$\pi_\la(f) = f - (1 - \cF_+)(1-F_\la)^{-1}P_\la(f)$$
%%
\end{corollaire}

\section{Deux remarquables distributions ayant la propriŽtŽ de Sonine}

Nous appliquons formellement la projection orthogonale aux
distributions paires invariante et anti-invariante sous Fourier:
$\delta_{-\la}(t) + \delta_\la(t) + 2\cos(2\pi\la t)$ et
$\delta_{-\la}(t) + \delta_\la(t) - 2\cos(2\pi\la t)$ et cela nous
mne ˆ 

\begin{definition}
On pose
%%
$$A_\la(t) = \delta_{-\la}(t) + \delta_\la(t) + 2\cos(2\pi\la t) -
\Big((1 + \cF_+)(1+F_\la)^{-1}(2\cos(2\pi\la y))\Big)(t)$$
$$-iB_\la(t) = \delta_{-\la}(t) + \delta_\la(t) - 2\cos(2\pi\la t) +
\Big((1 - \cF_+)(1-F_\la)^{-1}(2\cos(2\pi\la y))\Big)(t)$$
\end{definition}

\begin{theoreme}
Les distributions paires tempŽrŽes $A_\la(t)$ et $-iB_\la(t)$ ont la
propriŽtŽ de Sonine: elles sont nulles et de Fourier nulles dans
l'intervalle $(-\la,\la)$. Elles sont respectivement invariante et
anti-invariante sous Fourier.
\end{theoreme}

\begin{definition}
On pose
%%
$$\psi_+^\la(t) = 2\cos(2\pi\la t) -
\cF_+\Big((1+F_\la)^{-1}P_\la(2\cos(2\pi\la y))\Big)(t)$$
%%
$$\psi_-^\la(t) = 2\cos(2\pi\la t) +
\cF_+\Big((1-F_\la)^{-1}P_\la(2\cos(2\pi\la y))\Big)(t)$$
\end{definition}

\begin{theoreme}
La fonction entire paire $\psi_\pm^\la$ est (l'unique) fonction
entire dont la restriction ˆ $(-\la,\la)$ est $(1\pm
F_\la)^{-1}(2\cos(2\pi\la y))$. Elle est aussi \break{l'unique} fonction
solution de l'Žquation $\phi(x) \pm \cF(\phi|_{(-\la,\la)})(x) =
2\cos(2\pi \la x)$. On a $A_\la = \psi_+^\la + \cF(\psi_+^\la)$ et
$-iB_\la = - \psi_-^\la + \cF(\psi_-^\la)$. 
\end{theoreme}

\section{Les fonctions $\cE(w)$, $\cA(w)$, $\cB(w)$ pour les espaces
de Sonine}

Nous avons dŽmontrŽ le rŽsultat suivant:

\begin{theoreme}
Soit 
$$\cE_\la(w) = \pi^{-w/2}\Gamma(w/2)\Big(\la^{1/2 - w} +
\frac{\sqrt{\la}}{2}\int_\la^\infty (\psi_+^\la(t) -
\psi_-^\la(t))t^{-w}dt\Big)$$
L'intŽgrale est absolument convergente pour $\Re(w)>0$. La fonction
$\cE_\la(w)$ est une fonction entire satisfaisant la condition de
de~Branges pour le demi-plan $\Re(s)>1/2$ et l'espace $\cB(\cE_\la)$
co•ncide (isomŽtriquement) avec l'espace $\cS_\la$ des transformŽes de
Mellin complŽtŽes des fonctions de $K_\la$.
\end{theoreme}

Nous esquissons la dŽmonstration. On notera $X_w^\la(t)$ l'unique
ŽlŽment de $K_\la$ tel que $\forall f\in K_\la\ \int_0^\infty
f(t)t^{-w}dt = \int_0^\infty f(t)X_w(t)dt$ (pour $\Re(w)\leq 1/2$ la
premire intŽgrale est un prolongement analytique). Pour $\Re(w)>1/2$
la formule pour la projection orthogonale permet d'Žcrire
``explicitement'' $X_w^\la(t)$, et on voit en particulier que
$X_w^\la(t)- t^{-w}$ est la restriction ˆ $t>\la$ d'une fonction
entire. L'idŽe est de calculer la distribution $(td/dt +
w)X_w^\la(t)$: en effet on peut montrer en partant de la formule
\ref{noyau} que (la transformŽe de Mellin complŽtŽe de) la partie
invariante sous Fourier de cette distribution sera un multiple (rŽel
positif lorsque $w\in(1/2,+\infty)$) d'une fonction $\cA_\la(z)$
possible, et que (la transformŽe de Mellin complŽtŽe de) la partie
anti-invariante sous Fourier sera un multiple d'une fonction
$-i\cB_\la(z)$.  Or le calcul montre que la partie invariante de
$(td/dt + w)X_w^\la(t)$ est multiple de la distribution remarquable
$A_\la(t)$ dŽfinie prŽcŽdemment et que la partie anti-invariante est
multiple de $-iB_\la(t)$. Une Žtude plus poussŽe permet d'affirmer que
la combinaison $\sqrt{\la}/2 (A_\la - iB_\la)(t)$ a comme transformŽe
de Mellin (complŽtŽe par le facteur Gamma) une fonction $\cE_\la(w)$
possible pour l'espace de Hilbert de de~Branges Sonine $\cS_\la$. La
dŽmonstration donne aussi:

\begin{theoreme}
Il y a co•ncidence entre $\cE_\la(w)$ et $\sqrt{\la}$ fois la valeur
du ``saut'' de l'Žvaluateur (``euclidien'') $Z_w^\la(t) =
\pi^{-w/2}\Gamma(w/2)X_w^\la(t)$ en $t= \la$.
\end{theoreme}

Il est intŽressant que certaines quantitŽs apparaissant dans cette
Žtude et liŽes au noyau de Dirichlet jouent un r™le dans l'Žtude de
son dŽterminant de Fredholm sur un intervalle fini, dŽterminant dont
on conna”t l'importance dans la thŽorie des matrices alŽatoires
\cite{mehta}.



%\vfill\eject

\bibliographystyle{amsplain}
\begin{thebibliography}{10}


\bibitem{bra64} L.~de~Branges, \emph{Self-reciprocal
functions}, J. Math. Anal. Appl. {\bf 9} (1964) 433--457.

\bibitem{bra} L.~de~Branges, \emph{Espaces de Hilbert de Fonctions Entires},
Masson, Paris, 1972.

\bibitem{jfcras} J.-F.~Burnol, \emph{Sur certains espaces de Hilbert de
fonctions enti\`eres{,} li\'es \`a la transformation de Fourier et aux
fonctions L de Dirichlet et de Riemann},  C. R. Acad. Sci. Paris {\bf
333} (2001), s\'erie~I, 201-206.

\bibitem{jfhab}
J.-F. Burnol, \textit{On Fourier and Zeta(s)}, Habilitationsschrift,
\texttt{math/0112254}. Article soumis.

\bibitem{jfcopoisson}
J.-F. Burnol, \textit{Co-Poisson Intertwining: Distribution and Function Theoretic Aspects}, en prŽparation.


\bibitem{dymkean1}
H.~Dym, H.P. McKean, \textit{Fourier series and integrals}, Academic Press,
1972.

\bibitem{dymkean2}
H.~Dym, H.P.~McKean, \emph{Gaussian processes, function
theory, and the inverse spectral problem}, Academic Press, New York-London, 1976.

\bibitem{mehta}
M.~L.~Mehta, \emph{Random Matrices}, 2nd ed., Academic Press, San Diego, 1991.

\bibitem{slepian}
D.~Slepian, \emph{Some comments on Fourier analysis, uncertainty and modeling}, SIAM Review {\bf 25} (1983), No. 3, 379-393.

%\bibitem{krein} M.G.~Krein, \textit{Theory of entire functions of
%    exponential type} (in Russian), Izv. Akad. Nauk. SSSR,
%  Ser. Mat. 11 (1947), No. 4, 309-326.

%\bibitem{pal} R.E.A.C. Paley, N. Wiener,
%\textit{Fourier Transforms in the Complex Domain},
%Amer. Math. Soc., Providence, Rhode Island, 1934.
%
%\bibitem{tit}
%E. C. Titchmarsh, \textit{The Theory of the Riemann-Zeta Function},
%Clarendon Press, Oxford, 1951.

\end{thebibliography}

\end{document}
